% !PRINTED 297
\weberpage{297}
Nach dem Begriff des Normalteilers enthält $\mathfrak{N}$ auch die Substitution
\weq{12}{\begin{pmatrix} 0, & -1 \\ 1, & 0 \end{pmatrix}\begin{pmatrix} 0, & 1 \\ -1, & \xi \end{pmatrix}\begin{pmatrix} 0, & 1 \\ -1, & 0 \end{pmatrix} = \begin{pmatrix} \xi, & 1 \\ -1, & 0 \end{pmatrix}}
und folglich auch
\weq{13}{\begin{pmatrix} 0, & 1 \\ -1, & \xi \end{pmatrix}\begin{pmatrix} \xi, & 1 \\ -1, & 0 \end{pmatrix} = \begin{pmatrix} 1, & 0 \\ 2\xi, & 1 \end{pmatrix} \equiv A^{2\xi}\ [\S\ 79,\ (11)],}
also auch $A$ und alle seine Potenzen. Daher enthält $\mathfrak{N}$ auch
\weq{14}{\begin{pmatrix} 1, & 0 \\ \eta, & 1 \end{pmatrix}\begin{pmatrix} 0, & 1 \\ -1, & \xi \end{pmatrix} = \begin{pmatrix} 0, & 1 \\ -1, & \xi+\eta \end{pmatrix}}
für ein beliebiges $\eta$, d.\,h. die Substitution (11) für \gesperrt{jedes beliebige} $\xi$ und mithin auch
\[
B = \begin{pmatrix} 0, & 1 \\ -1, & 0 \end{pmatrix}.
\]
Demnach ist nach dem Satze des vorigen Paragraphen die Gruppe $\mathfrak{N}$ mit $\mathfrak{L}_0$ identisch.

Es bleibt noch die Möglichkeit zu erörtern, daß in (11) $\xi \equiv 0$ ist.

In diesem Falle enthält die Gruppe $\mathfrak{N}$ also die Substitution
\[
B = \begin{pmatrix} 0, & 1 \\ -1, & 0 \end{pmatrix},
\]
und folglich auch für ein beliebiges, durch $p$ nicht teilbares $\alpha$
\[
V = \begin{pmatrix} 0, & \alpha \\ -\alpha^{-1}, & 0 \end{pmatrix}\begin{pmatrix} 0, & 1 \\ -1, & 0 \end{pmatrix}\begin{pmatrix} 0, & -\alpha \\ \alpha^{-1}, & 0 \end{pmatrix}\begin{pmatrix} 0, & -1 \\ 1, & 0 \end{pmatrix} = \begin{pmatrix} \alpha^2, & 0 \\ 0, & \alpha^{-2} \end{pmatrix}.
\]
Daraus leitet man, als in $\mathfrak{N}$ enthalten, noch weiter ab:
\[
W = \begin{pmatrix} 1, & 0 \\ -1, & 1 \end{pmatrix}\begin{pmatrix} \alpha^{-2}, & 0 \\ 0, & \alpha^2 \end{pmatrix}\begin{pmatrix} 1, & 0 \\ 1, & 1 \end{pmatrix}\begin{pmatrix} \alpha^2, & 0 \\ 0, & \alpha^{-2} \end{pmatrix} = \begin{pmatrix} 1, & 0 \\ \alpha^4-1, & 1 \end{pmatrix}.
\]
Wenn nun $p$ größer ist als 5, so kann man $\alpha$ so annehmen, daß $\alpha^4 - 1$ nicht durch $p$ teilbar ist; dann enthält also $\mathfrak{N}$ auch die Substitution $A$ und ihre Potenzen und mithin ist $\mathfrak{N}$ mit $\mathfrak{L}_0$ identisch.

Ist $p = 5$, so ist $\alpha^4 - 1$ immer durch $p$ teilbar, also dieser Schluß nicht anwendbar.

In diesem Falle folgt aber aus dem Begriff des Normalteilers als in $\mathfrak{N}$ enthalten:
%CONT>>
